% Generated draft: requires mathematical and source review.
\documentclass[11pt]{article}
\usepackage[margin=25mm]{geometry}
\usepackage{fontspec}
\setmainfont{Latin Modern Roman}
\usepackage{amsmath,amssymb,mathtools,unicode-math}
\setmathfont{Latin Modern Math}
\usepackage{enumitem,needspace,xurl,hyperref}
\setlength{\parindent}{0pt}
\setlength{\parskip}{5pt}
\setlength{\emergencystretch}{3em}
\newcommand{\N}{\mathbb{N}}
\newcommand{\Z}{\mathbb{Z}}
\newcommand{\Q}{\mathbb{Q}}
\newcommand{\R}{\mathbb{R}}
\newcommand{\C}{\mathbb{C}}
\newcommand{\F}{\mathbb{F}}
\newcommand{\A}{\mathbb{A}}
\newcommand{\PP}{\mathbb{P}}
\newcommand{\E}{\mathbb{E}}
\newcommand{\eps}{\varepsilon}
\newcommand{\dd}{\,\mathrm{d}}
\begin{document}
\section*{UnsolvedMath ID: 30001702}
\textbf{Source problem code:} catalogue entry
\subsection{Definitions and mathematical statement}
Let \(S^1=\{z\in\C:|z|=1\}\), and let \(T^d=(S^1)^d\) denote the \(d\)-dimensional torus. Under the usual regular-cell interpretation, a simplicial cell decomposition \(P\) is a regular CW decomposition whose closed cells, with their face structures, are combinatorially simplices. Distinct cells may have the same vertex set; thus this interpretation does not require \(P\) to be an abstract simplicial complex. A facet is a maximal cell; in a pure \(d\)-dimensional decomposition these are precisely the \(d\)-cells. Write \(f_d(P)\) for the number of \(d\)-cells.

The supplied question is whether, for every dimension \(d\) and every simplicial cell decomposition \(P\) of \(T^d\),
\[
 f_d(P)\geq (d+1)!.
\]
The source record does not explicitly specify the range of \(d\), finiteness conventions, or the precise meaning of ``simplicial cell decomposition.'' The regular-cell interpretation above therefore requires confirmation against the cited report. Dimensions are understood to be nonnegative integers; whether the source intends to include \(d=0\), with \(T^0\) a point, is unspecified.
\subsection{Short English statement}
Must every simplicial cell decomposition of the \(d\)-dimensional torus have at least \((d+1)!\) facets? The supplied record labels this question open and reports a Steinberg torus construction attaining that number. These are supplied source assessments, and the exact cell-decomposition conventions require review.
\subsection{Sources}
\begin{itemize}
\item Local UnsolvedMath registry, permanent ID \(30001702\), title: \emph{Minimal Facets in Triangulated Tori}; source problem code \texttt{catalogue entry}, in the Oberwolfach Reports Open Problems collection. The record assigns the year 2011 and supplies no individual proposer.
\item \emph{Topological and Geometric Combinatorics} (2011), Oberwolfach Reports, pp.\ 349--423, DOI \texttt{10.4171/owr/2011/08}. Supplied source URL and literature reference labeled ``original Oberwolfach report'': \url{https://doi.org/10.4171/owr/2011/08}.
\item Supplied literature assessment, checked on \texttt{2026-08-21}: the question is labeled open; the Steinberg torus is reported to provide \((d+1)!\) facets, while the matching lower bound is described as an open problem for which the supplied review found no general proof. This is supplied dated curation, not an independently verified literature assessment.
\item The registry labels the statement \texttt{corrected\_verified} and reports a manual semantic, typographic, and LaTeX review. That supplied editorial qualification does not establish mathematical correctness or verify the current research status.
\end{itemize}
\end{document}
